\documentclass[leqno, a4paper, 12pt]{jsarticle}
\usepackage{amssymb}
\usepackage{amsthm}
\usepackage{amsmath}
\usepackage{comment}
	%可換図式用
		%\usepackage[all]{xy}
	%重くなるので使わいときはコメント化しておく。
%定理環境 thmはあまり使わずpropをなるべく使う。
%主張は基本的に証明の中で使い、そのため番号を付けない。
%注意環境を付けてみた。
\newtheorem{thm}{定理}
\newtheorem{prop}[thm]{命題}
\newtheorem{cor}[thm]{系}
\newtheorem{lem}[thm]{補題}
\newtheorem{df}[thm]{定義}
\newtheorem{exam}[thm]{例}
\newtheorem{fact}[thm]{事実}
\newtheorem*{claim}{主張}
\newtheorem*{rmk}{注意}
\renewcommand{\proofname}{{\bf 証明}}
%矢印たち。
%\toは\rightarrowと同じ。\getsは\leftarrowと同じ。同値は\iff
%上向き矢はuparrow逆はdownarrow
\newcommand{\To}{\Longrightarrow}
\newcommand{\Gets}{\Longleftarrow}
%黒板太字
\newcommand{\nn}{\mathbb{N}}
\newcommand{\zz}{\mathbb{Z}}
\newcommand{\qq}{\mathbb{Q}}
\newcommand{\rr}{\mathbb{R}}
\newcommand{\cc}{\mathbb{C}}
%無限大
\newcommand{\mgn}{\infty}
%ギリシャン
\newcommand{\dpst}{\displaystyle}
\newcommand{\ep}{\varepsilon}
\newcommand{\al}{\alpha}
\newcommand{\be}{\beta}
\newcommand{\de}{\delta}
\newcommand{\la}{\lambda}
\newcommand{\La}{\Lambda}
\newcommand{\ga}{\gamma}
\newcommand{\lil}{\la\in \La}
%商集合
\newcommand{\qt}[2]{#1/\! #2}
%enumerate環境
\renewcommand{\labelenumi}{{\rm (\arabic{enumi})}}
%以降alignじゃなくてenumerate を使え。
%なんか演算
\DeclareMathOperator{\card}{card}
\DeclareMathOperator{\supp}{supp}
%なんか演算２
\DeclareMathOperator{\bd}{Bdry}
\DeclareMathOperator{\cl}{CL}
\DeclareMathOperator{\nai}{Int}
\DeclareMathOperator{\di}{diam}

%差集合は\setminusで統一する。等号の否定は\neq
%真部分集合は\subsetneqq　偏微分は\partial
%ディスプレイスタイル。\dfracでディスプレイ分数
%空集合は\Oを使う！！！！！！！！！！！！

\title{バナッハ空間の関数空間への埋め込みと，関数空間の可分距離空間に対する普遍性}
\author{ナンブ　キトラ}
\date{\today}
			\begin{document}
\maketitle

この文章では，まず最初に任意の実バナッハ空間がコンパクト空間上の実連続関数の空間に等長に埋め込めることを示す．次に任意の可分距離空間を等長的に埋め込めるような関数空間を調べる．この文章では実係数の線形空間しか扱わないので，バナッハ空間と言ったら全て実係数である．

\section{バナッハ空間の関数空間への埋め込み}


一般にバナッハ空間$B$に対して$B^*$でその双対空間を表す．
また，
$f\in B^*$と$x\in B$に対して
$\langle f, x\rangle=f(x)$と定義する．
そして$B^*$上のノルム$\|*\|_{B^*}$は
\[
\|f\|_{B^*}=\sup_{x\in B, x\neq 0}\frac{|f(x)|}{\|x\|_B}=
\sup_{\|x\|_B\le 1, x\neq 0}\frac{|f(x)|}{\|x\|_B}
=
\sup_{\|x\|_B\le 1}|f(x)|
=\sup_{\|x\|_B=1}|f(x)|.
\]
と定義する．


一般にバナッハ空間$B$のノルムは単に$\|*\|$と書くが，区別の必要があるときには$\|*\|_B$と書く．


以下の定理はHahn--Banachの定理の系の一つである．Hahn--Banachの定理については例えば\cite{kaiseki}を
参照のこと．
\begin{thm}[Hahn--Banach]
$B$を$0$次元線形空間でないバナッハ空間とし，$x\in B$とする．
このとき，$f\in B^*$で，
$\|f\|_{B^*}=1$で$f(x)=\|x\|_B$
となるものが存在する．
\end{thm}


\begin{cor}\label{cor:norm}
任意のバナッハ空間$B$と任意の$x\in B$について
\begin{align*}
\|x\|_B=\sup_{f\in B^*, f\neq 0}\frac{|f(x)|}{\|f\|_{B^*}}=
\sup_{\|f\|_{B^*}\le 1, f\neq 0}\frac{|f(x)|}{\|f\|_{B^*}}=
\sup_{\|f\|_{B^*}\le 1}|f(x)|
=
\sup_{\|f\|_{B^*}=1}|f(x)|.
\end{align*}
が成り立つ．
\end{cor}

$B$をバナッハ空間とする．$S\subset B^*$とする．
$x\in B$に対して$\tau_x(f)=\langle f, x\rangle $として写像$\tau_x: B^*\to \rr$を定義し，写像の族
$\{\tau_x|_S: S\to \rr\}_{x\in B}$の元の全てを連続にする$S$上の位相のことを
\textbf{$S$上の汎弱位相}と呼ぶ．
$S$上の汎弱位相は$B^*$上の汎弱位相を$S$に制限した位相に等しい．
また，$S$上の汎弱位相はハウスドルフ位相である．
以下では，$\{f\in B^*\mid \|f\|\le 1\}$に汎弱位相を入れた空間を，$U[B]$と書くことにする．

次の有名な定理の証明は，例えば\cite{kaiseki}を参照されたい．
\begin{thm}[Banach--Alaoglu]
$B$をバナッハ空間とする．このとき汎弱位相を持つ空間
$U[B]=\{f\in B^*\mid \|f\|\le 1\}$
はコンパクト空間である．
\end{thm}

以上のBanach--Alaogluの定理を使って，バナッハ空間を関数空間に埋め込む．その前に関数空間を定義する．




\begin{df}
空でない位相空間$X$に対して$C_b(X)$で$X$から$\rr$への有界実連続関数全体を表すことにする．
そして$f, g\in C_b(X)$に対して，
\[
\|f-g\|_{C_b(X)}=\sup_{x\in X}|f(x)-g(x)|
\]
で$C_b(X)$上のノルム$\|*\|_{C(X)}$
を定義する．
このノルムを単に$\|*\|$と省略したりする．
このとき$C_b(X)$はバナッハ空間になる．

同様に$C(X)$で$X$上の実連続関数
全体の集合とし，上と同様に$\sup$ノルムを導入する．このときノルムは$\infty$の値を一般にとりうる．

$X$がコンパクト空間の場合には
$C(X)=C_b(X)$となることに注意せよ．
\end{df}

以下の定理がこのセクションの主定理である．
\begin{thm}\label{thm:CK}
任意のバナッハ空間$B$に対して，
$\tau: B\to C(U[B])$を任意の$x\in B$に対して，
$\tau(x)(f)=\langle f, x\rangle $で定義する．
このとき，$\tau$は線形等長写像である．
よって，特に任意のバナッハ空間はコンパクトハウスドルフ空間上の
関数空間に等長埋め込むことができる．
\end{thm}
\begin{proof}
任意の$x\in B$について$\tau(x)$が$U[B]$上の連続関数であることは，
汎弱位相の定義そのものから従う．
$\tau$の線形性は明らかである．
等長性を示そう．
任意の$x\in B$について，
\[
\|\tau(x)\|_{C(U[B])}=\|x\|_B
\]
を示せば良い．

系\ref{cor:norm}より，

\[
\|\tau(x)\|_{C(U[B])}=\sup_{f\in U[B]}|\tau(x)(f)|
=\sup_{\|f\|_{B^*}\le 1}|f(x)|=\|x\|_B
\]
なので，証明は終わる．
\end{proof}


%%%%%%%%%%%%%%%%%%%%%%%%%%%%%%%%%%%%
%%%%%%%%%%%%%%%%%%%%%%%%%%%%%%%%%%%%
%%%%%%%%%%%%%%%%%%%%%%%%%%%%%%%%%%%%
%%%%%%%%%%%%%%%%%%%%%%%%%%%%%%%%%%%%
%%%%%%%%%%%%%%%%%%%%%%%%%%%%%%%%%%%%
%%%%%%%%%%%%%%%%%%%%%%%%%%%%%%%%%%%%
%%%%%%%%%%%%%%%%%%%%%%%%%%%%%%%%%%%%
%%%%%%%%%%%%%%%%%%%%%%%%%%%%%%%%%%%%

\section{関数空間の可分距離空間に対する普遍性}
\subsection{普遍性の定理}
このセクションでは，任意の可分距離空間を等長に埋め込むことができるような関数空間を調べる．

二つの空でない位相空間$X$, $Y$と連続写像$p: X\to Y$に対して，
$C(p): C(Y)\to C(X)$を$C(p)(f)=f\circ p$で定義する．
このとき以下が成り立つ．
\begin{prop}\label{prop:hannpenn}
二つの空でない位相空間$X$, $Y$と連続写像$p: X\to Y$に対して，
$p$が全射ならば，$C(p): C(Y)\to C(X)$はノルムを保つ．つまり
任意の$f\in C(Y)$について，
\[
\|f\|=\|C(p)(f)\|
\]
となる．ここで，ノルムには$\infty$の値も許されていることに注意しよう．
\end{prop}
\begin{proof}
任意の$f\in C(Y)$について，$p$の全射性から
\[
\|f\|=\sup_{y\in Y}|f(x)|=\sup_{x\in X} |f\circ p(x)|=\|f\circ p\|
\]
なので，証明が終わる．
\end{proof}

空でない距離空間$X$に対して，$d$をその距離とし，$o\in X$を固定点とする．このとき$x\in X$に対して
$K(x): X\to \rr$を
\[
K(x)(y)=d(x, y)-d(y, o)
\]
と定義することによって
写像$K: X\to C_b(X)$
を定義する．このとき
$K$は等長埋め込みになっている．
この埋め込み写像をクラトフスキー埋め込みと呼ぶ．

\begin{lem}\label{cor:banach}
可分距離空間$X$に対して，
可分バナッハ空間$B$が存在して，
$X$から$B$への等長埋め込みが存在する．
\end{lem}
\begin{proof}
$K: X\to C_b(X)$をクラトフスキー埋め込みとする．
このとき$K(X)$によって生成される$C_b(X)$の部分空間の完備化が求める可分バナッハ空間になる．
\end{proof}

\begin{df}
距離空間$X$が
\textbf{可分距離空間に対して普遍的}
であるとは，任意の可分距離空間$A$
に対して，$A$から$X$への等長埋め込みが存在するときにいう．
\end{df}


\begin{df}
位相空間$X$が\textbf{パラメトライジング}であるとは，
任意の可分バナッハ空間$B$と汎弱位相を持つ双対空間の単位球$U[B]=\{f\in B^*\mid \|f\|\le 1\}$
に対して，$X$から$U[B]$に対して全射連続写像が存在するときにいう．
この用語はこの文章でしか使わないので注意されたい．
\end{df}


以下の定理がこの文書の中心的な定理である．証明は\cite[Theorem 3.6]{H2}の証明を大いに参考にした．
\begin{thm}
$X$をパラメトライジングな位相空間とする．
このとき，$C_b(X)$は可分距離空間に対して普遍的である．
\end{thm}
\begin{proof}
$T$を可分距離空間とする．
このとき$C_b(X)$に$T$
が埋め込めることを証明する．
系\ref{cor:banach}
により，可分バナッハ空間$B$と等長埋め込み
$h: T\to B$が存在する．
さて，
$p: X\to U[B]$を全射連続写像とする．
$X$がパラメトライジングであることからこのような写像は存在する．
このとき$p: X\to U[B]$
は全射なので，命題\ref{prop:hannpenn}から
$C(p): C(U[B])\to C_b(X)$
は等長埋め込みである．
また，定理\ref{thm:CK}から等長埋め込み
$\tau: T\to C(U[B])$が存在する．よって
写像$C(p)\circ \tau\circ h: T\to C_b(X)$が求める等長埋め込みである．
\end{proof}

\subsection{普遍性を持つ関数空間を持つ空間}
このサブセクションでは可分距離空間に対して普遍性を持つ距離空間の例を挙げていく．最初にパラメトライジングなもの，後にそうでないものを挙げる．

以下の定理は基本的なものである．
\begin{lem}\label{lem:kabunn}
$B$を可分バナッハ空間として汎弱位相を持つ空間$U[B]=\{f\in B^*\mid \|f\|\le 1\}$はコンパクトで，距離化可能空間であり，
弧状連結で，かつ$U[B]$の任意の空でない開集合も弧状連結である．
\end{lem}

以下の定理は結構有名である．証明は例えば
\cite[Theorem 4.18]{K}を参照のこと．

\begin{thm}
$\Gamma$をカントール集合とする．
このとき，
任意の空でないコンパクト距離化可能空間 $K$
に対して全射連続写像$p: \Gamma\to K$が存在する．
特に$\Gamma$
はパラメトライジングである．
\end{thm}

以下の定理の証明は，例えば\cite{yamyam}を参照のこと．
\begin{thm}[Hahn--Mazurkiewicz]
空でない位相空間$X$
は以下を満たすとする．
\begin{enumerate}
\item $X$はコンパクトな距離化可能空間
\item $X$は連結
\item $X$は局所連結，つまり，連結開集合からなる開基が存在する．
\end{enumerate}
このとき，全射連続写像$p: [0, 1]\to X$
が存在する．
特に$[0, 1]$はパラメトライジングである．
\end{thm}

距離化可能空間$X$が\textbf{絶対レトラクト}であるとは，
$X$が他の距離空間$Y$に閉集合として位相的に埋め込まれているならば，
連続写像$R: Y\to X$が存在して，$R|_{X}=id_X$を満たす．ここで，$X\subset Y$とみなしている．

次の定理は単位閉区間$[0, 1]$の可算直積空間$[0, 1]^{\nn}$の特徴付け定理の一つである．空間$[0, 1]^{\nn}$は\textbf{ヒルベルトキューブ}と呼ばれる．
以下の定理は\cite{Keller}によってフレシェ空間の場合に原型が示され，
\cite{0}によって以下に述べる形になった．
\begin{thm}[Keller-Dobrowolski--Toru\'nczyk]\label{KDT}
$K$を位相線形空間の凸集合とする．
もし，$K$は無限大の位相次元を持ち，コンパクトで距離化可能空間であって，絶対レトラクトならば，
$K$は$[0, 1]^{\nn}$に同相
である．
\end{thm}

距離空間が絶対レトラクトかどうかを証明するには拡張定理がとても役に立つ．
次の定理はDugundjiの拡張定理と呼ばれている．証明はこのブログでも紹介している \cite{con}．また原論文\cite{1} を参照のこと．
局所凸空間の定義は\cite{kaiseki}などを参照のこと．大雑把に局所凸位相線形空間のことを説明すると，セミノルムの族によって位相が生成されている位相線形空間が局所凸空間である．
\begin{thm}[Dugundji]\label{d}
距離空間$(X,d)$の閉集合$A$と局所凸型線形位相空間$L$の凸集合$C$及び連続写像$f:A\to C$について、連続写像$F:X\to C$が存在し、
\[
F|A=f
\]
を充たす。
\end{thm}

これを用いると次のことがわかる．
\begin{prop}
任意の無限次元可分バナッハ空間$B$について，$U[B]$は絶対レトラクトである．
特に$U[B]$はヒルベルトキューブ$[0, 1]^{\nn}$に同相である．
\end{prop}
\begin{proof}
 $U[B]$が他の距離化可能空間$Y$に閉集合として位相的に埋め込まれているとする．
 ここでは$U[B]\subset Y$とみなすことにする．
$B^*$に汎弱位相を導入すると，この位相空間は局所凸線形位相空間になるので，
 恒等写像$id: U[B]\to U[B]$は，
Dugundjiの拡張定理(定理\ref{d})の仮定を満たす．
したがって$id$を$Y$全体に拡張し，連続写像$R: Y\to U[B]$であって，
 $R|_{U[B]}=id$なるものを得る．よって$U[B]$は絶対レトラクトである．
 補題\ref{lem:kabunn}より，
 $B$が可分なので，$U[B]$はコンパクトで距離化可能である．ゆえに定理\ref{KDT}から，$U[B]$がヒルベルトキューブに同相であることがわかる．
\end{proof}

有限次元バナッハ空間$B$について$U[B]$は有限次元閉球に同相で，この空間は
$[0, 1]^{\nn}$からの全射連続写像が存在するので，以下を得る．
\begin{cor}
ヒルベルトキューブ$[0, 1]^{\nn}$はパラメトライジングである．
\end{cor}

以上であげた三つの空間がパラメトライジングであることは以下のようにまとめることができる．
\begin{cor}
バナッハ空間
$C(\Gamma)$，$C([0, 1])$そして$C([0, 1]^{\nn})$は可分距離空間に対して普遍的である．
\end{cor}

\begin{rmk}
$C(\Gamma)$と$C([0, 1])$が可分距離空間に対して普遍性を持つことは，
\emph{\cite{B}}で最初に示された．
\end{rmk}




一般に距離空間$X$がパラメトライジングでなくとも，その上の関数空間は可分距離空間に対して普遍的であり得る．
無限次元可分バナッハ空間$B$について$U[B]$は非可算集合なので，
可算空間はパラメトライジングではないことに注意しよう．
\begin{prop}
有理数の空間$\qq$について，
$C_b(\qq)$は可分距離空間に対して
普遍性を持つ．
\end{prop}
\begin{proof}
順序を保つ位相的埋め込み
$H: \qq\to [0, 1]$で$H(\qq)$が$[0, 1]$のなかで稠密になるものを考える．
任意の可分距離空間$X$と
等長埋め込み$I: X\to C([0, 1])$
をとる．
このとき，
写像$J: X\to C_b(\qq)$
を
$J(x)=I(x)\circ H$と定義する．
すると，$H(\qq)$が$[0, 1]$
の中で稠密であることから，
$J$が等長埋め込みであることがわかる．
\end{proof}

\begin{prop}
$N$は可算離散空間とする．
このとき， $C_b(N)$
は可分距離空間に対して普遍的である．
\end{prop}
\begin{proof}
空間$N$について
\[
N=\{(k, n)\in \nn^2\mid n\in \nn, 0\le k\le 2^n\}
\]
と仮定しても良い．

任意の可分距離空間$X$と
等長埋め込み$I: X\to C([0, 1])$
をとる．
このとき，
写像$J: X\to C_b(N)$
を
\[
J(x)(k, n)=f\left(\frac{k}{2^n}\right)
\]
と定義する．
このとき，
$\{k/2^n\mid (k, n)\in N\}$
は$[0, 1]$の中で稠密であることから，
$J$は等長埋め込みになっている．
\end{proof}

カントール集合上の関数空間$C(\Gamma)$
が可分距離空間に対して普遍的であるということからより広いクラスの空間に対して同様の性質が言える．


以下の定理はDugundjiの拡張定理の証明法を詳細に見るとわかることである．\cite[Theorem 5.1]{1}を参照のこと．
\begin{prop}\label{prop:d}
$X$を空でない距離化可能空間とし，$A$をその空でない閉集合とする．
このとき，
等長な線形写像
\[
\phi: C(A)\to C(X)
\]
が存在し，任意の$f\in C(A)$に対して
$\phi(f)|_A=f$を満たす．
\end{prop}


\begin{prop}
$X$を非可算可分完備距離化可能空間とする．
このとき，
$C_b(X)$は可分距離空間に対して普遍的である．
\end{prop}
\begin{proof}
$X$が非可算可分完備距離化可能空間ならば，$X$はカントール集合
$\Gamma$
と同相な部分集合を持つ(例えば\cite[Corollary 6.5]{K}を参照のこと)．
よって，命題\ref{prop:d}
から，$C_b(X)$は$C(\Gamma)$
と等長的な部分空間を含むので，
特に可分距離空間に対して普遍的である．
\end{proof}




\begin{rmk}
$C(\omega_0+1)$
は可分距離空間に対して普遍的でない．
特に，$C(\omega_0\cdot 2+1)$
は$C(\omega_0+1)$
に等長的に埋め込めない．
これは\emph{\cite[Lemma 5]{V}}を参照のこと．
\end{rmk}













	
\begin{thebibliography}{99}
\bibitem{kaiseki}宮島静雄, 関数解析, 
横浜図書, 2005年.
\bibitem{yamyam}yamyamtop, 全射な曲線と単射な曲線, \\
https://yamyamtopo.files.wordpress.com/2019/08/hahn-mazurkiewicz.pdf, 
2020年9月閲覧．
\bibitem{con}電波通信(ナンブキトラ), ダガンディーとハウスドルフの拡張定理, 
http://concious4410.hatenablog.com/entry/2017/01/18/185814

\bibitem{B}S. Banach, \textit{Th\'eorie des \'Operations Lin\'eaires},  
Chelsea Publishing Co., New
York, 1955. 
\bibitem{0}T. Dobrowolski and H. Tru\'nczyk, 
\textit{On metric linear space homeomorphic to $\ell^2$ and compact convex sets homeomorphic to $Q$}, Bull. Acad. Polon. Sci. S\'er. Sci. Math. 
27 (1979), 883--887. 
\bibitem{1}J. Dugundji, \textit{An extension of Tietze's theorem}, 
Paciffic J. Math. vol. 1 (1951), 105--125. 
\bibitem{H2}J. Heinonen, \textit{Geometrical embedding of metric spaces}, Report, University of Jyv\"askyl\"a Department of Mathematics and Statics, vol. 90. 2003. http://www.math.jyu.fi/tukimas/ber.html.
\bibitem{K}A. S. Kechris, \textit{Classical Descriptive Set Theory}, 
Springer-Verlag, GTM 156 (1995). 
\bibitem{Keller}O. H. Keller, \textit{Die homoiomorphie der kompakten konvexen mengen in Hirbertscen Raum}, 
Math. Ann. 105 (1931), 748--758. 
\bibitem{V}R. Villa, \textit{Isometric embedding into space of continuous functions}, 
Studia Math. 129 (3) (1998), 
197--205. 
\end{thebibliography}




			\end{document}



\begin{comment}
[集合のやつは$\mid$を使う。]
[真なる包含関係]
\subsetneqq
[可換図式]
\[\xymatrix{
&   & (2) \ar@{=>}[rd]&  &  &  & (5) \ar@{=>}[rd]&\\ 
& (1)\ar@{=>}[ru] \ar@{=>}[rd]&  &(4)\ar@{=>}[r]&(7)& (1)\ar@{=>}[ru] \ar@{=>}[rd]& & (7)&\\ 
&   & (3) \ar@{=>}[ur]&  &  &  &(6) \ar@{=>}[ur]&
}\]

[\bigin \endを使わない方法]

{\prop[aa] aaa}
で
\begin{prop}[aa]
aaa
\end{prop}
と同じ\df \thm \proof でも同様

[直和]
$\amalg$

[同値関係でよく使う波線]
\sim

[引数付きの新命令]
\newcommand{\prn}[1]{\|#1\|_p}

[番号のリセット]

\setcounter{equation}{0}


[字体の変更]

{\rm hogehoge}と使う。

\rm ローマン
\it イタリック
\sl 斜体

[supp]

\newcommand{\supp}{\text{{\rm supp}}}

[中央揃えの改行]

	\begin{eqnarray*}
		hoge1&=&hogehoge
		hoge2&=&hogehofe

	\end{eqnarray*}

&&の部分で揃えられる。


[\tag]
\begin{equation}
f(x) = ax^2 + bx + c \tag{1.2.3}
\end{equation}



[場合分け]

	\begin{cases}
		hoge1 & \text{状況１}\\
		hoge2 & \text{状況２}\\
		・
		・
		・
	\end{cases}



\end{comment}





